\documentclass[a4paper,12pt]{article}
\usepackage{color}
\usepackage{graphicx, gensymb}
\usepackage{amsfonts, cancel, amssymb}
\usepackage{amsmath, amsthm, array, esvect, esint}
\usepackage{typearea}
\usepackage[a4paper,left=2cm,right=2cm,top=2cm,bottom=3cm,bindingoffset=5mm]{geometry}
\newcommand\numberthis{\addtocounter{equation}{1}\tag{\theequation}}
\newcommand{\bra}{}
\newcommand{\kom}{}
\newcommand{\brak}{}
\newcommand{\braket}{}
\newcommand{\intx}{}
\newcommand{\intr}{}
\newcommand{\kla}{}
\newcommand{\klam}{}
\newcommand{\klammer}{}
\newcommand{\oper}{}
\newcommand{\tecxt}{}

\begin{document}

\title{Franck-Hertz experiment}
\author{Joachim Schwardt}
\date{\today}
\maketitle

\section{Tasks}
\begin{enumerate}
\item The anode current is measured as a function of the acceleration voltage of a Hg-filled electron tube for three different vapor temperatures. 
\item The anode current is measured again for a Ne-filled electron tube at three different control voltages. 
\item The optical spectrum of the light emitted from the Ne-filled tube is measured and compared to the respective energy level scheme.
\item Plot of the maxima of the $I-U$ curves of the Hg-filled tube. 
\item The average excitation energies of Ne and Hg are identified via the maxima of the anode current. The frequency and wavelength of the EM-radiation emitted from the atoms after collision is calculated. 
\end{enumerate}
\section{Theoretical remarks}
According to Bohr's postulates of quantum mechanics, atoms can only exist in stationary states of definite energy, which follow a discrete series $E_1,\dots,E_n$. Hence, the emission or absorption of electromagnetic radiation is only possible for frequencies of the form $f=\frac{1}{h}|E_m-E_n|$, which corresponds to a transition of one discrete state to another. 
\subsection{The Franck-Hertz curve}
Picture a Hg-filled electron tube with an oxid-cathode $C$, heated by $U_H=6,3\,$V, where electrons are accelerated towards an anode $A$  with an acceleration volatge of $U_1=0\dots 60\,$V. Along their way to the mesh-shaped anode, the electrons can collide elastically or inelastically with Hg-atoms, depending on their energy. The inelastic collision is accompanied by a state transition of the atomic shell or an ionisation of the atom. The electron dispenses the energy needed for the transition. Further to the right of $A$, there is another electrode $S$, which is relative-negatively charged compared to $A$ with $U_2=2\,$V. If the electrons have enough kinetic energy left after the collision to overcome this potential barrier, they contribute to the current $I_A$. \\
An increase of $U_1$ will therefor lead to an increase of $I_A$, as the electrons have more kinetic energy left after an inelastic collision, up to a local maximum. However, once the minimal excitation energy has been reached, a further increase of the acceleration voltage leads to a higher probability of inelastic scattering, which leads to less electrons reaching $S$ and a slight decline in $I_A$ towards a local minimum. Additional kinetic energy helps more electrons to overcome the counter voltage, even after an inelastic collision. This process may repeat several times, with the maxima of $I_A$ corresponding to an integer multiple of the excitation energy. \\
Since this curve is not strictly periodic, the excitation energy $E_a$ is calculated from the mean of the maxima of the acceleration voltage, times the elementary charge $e$, i.e. $E_a=e\cdot\bra\langle {\Delta U_1^{\tecxt\text{max}}}\rangle$.
\subsection{Child's law}
Let $m_e$ be the electron mass, $\vv{v}_D$ the drift velocity, $e$ the elementary charge and $\vv{E}$ the electric field generated by the acceleration voltage. According to Drude's theory one can obtain the following equation of motion:
\begin{align*}
\dot{\vv{v}}_D + \frac{\vv{v}_D}{\tau}=-e\vv{E}\numberthis\label{Drude equation}
\end{align*} 
The mean collision time $\tau=\frac{\lambda}{v_\tecxt\text{therm}}$ is defined by the mean free path $\lambda$ and the mean of the absolute value of the electron gas velocity. It can be calculated from the Maxwell-Boltzmann distribution as follows:
\begin{align*}
v_\tecxt\text{therm}&=\sqrt{\frac{8k_BT}{\pi m_e}}\numberthis\label{v therm}
\end{align*}
Here $T$ is the temperature of the electron gas and $k_B$ the Boltzmann-constant. Consider the stationary case and define the mobility of the electrons as $\mu:=\frac{e\tau}{m_e}=\frac{e\lambda}{m_e v_\tecxt\text{therm}}$:
\begin{align*}
\vv{v}_D&=-\mu\vv{E}\numberthis\label{Drude mu equation}
\end{align*}
We can calulate the electric field from the 1-D Maxwell equation. Let $\varrho$ be the charge density and $j=v_D\varrho$ the current density:
\begin{align*}
E'(x)&=\frac{\varrho(x)}{\epsilon_0\epsilon_r}=\frac{j}{\epsilon_0\epsilon_r v_D}\overset{(\ref{Drude mu equation})}{=}-\frac{j}{\epsilon_0\epsilon_r\mu E}\numberthis\label{E field Maxwell}
\end{align*}
The continuity equation $j'+\partial_t\varrho=0$ demands that $j$ is a constant in the stationary case. Seperation of variables in Eq. (\ref{E field Maxwell}) and $E(x=0)=0$ then leads to:
\begin{align*}
\intx\int_{0}^{x}E\,\mathrm{d}E&=-\frac{j}{\epsilon_0\epsilon_r\mu}\intx\int_{0}^{x}\,\mathrm{d}x'\ \ \Rightarrow\ \ E(x)=\sqrt{\frac{-2jx}{\epsilon_0\epsilon_r\mu}}\numberthis\label{E field}
\end{align*}
We have $\varphi'(x)=-E(x)$ and $U(x=d)=U_B$, where $U_B$ is the applied acceleration voltage, which gives:
\begin{align*}
U(x)&=\varphi(x)-\varphi(0)=-\intx\int_{0}^{x}E(x')\,\mathrm{d}x'=-\sqrt{\frac{-8jx^3}{9\epsilon_0\epsilon_r\mu}}\numberthis\label{U of x}
\end{align*}
Accounting for a small current $I_0(U_B=0)$ and a contact voltage $U_K=U_1-U_B$ we obtain Child's law by evaluating (\ref{U of x}) at $x=d$ and $I_A=j A_K$ where $A_K$ is the anode area:
\begin{align*}
I_A(U_1)&=I_0+j(d)A_K=I_0+\frac{9\epsilon_0\epsilon_r e\lambda A_K\sqrt{\pi}}{8d^3\sqrt{8m_ek_BT}}\cdot\kla\left( {U_1-U_K} \right)^2\numberthis\label{Child law}
\end{align*}
\section{Experiment}
\subsection{Mercury measurement}
Parameters:\\
Hg-tube, heating voltage $U_H=6,3\,V$, retarding voltage $U_2=2,0\,$V, $\vartheta=(175,\,180,\,185)\,\degree$C. Each measurement was repeated twice to get a total of 9 curves. See Fig. \ref{Hg plot 175 degree} for one of these as a general example.\\
\begin{figure}[h!]
\centering
\includegraphics[scale=1]{Hg_T_175_plot.png}
\caption{Hg-tube, $\vartheta=175\,\degree$C}
\label{Hg plot 175 degree}
\end{figure}
Taking the average for the individual temperatures (superscript $\equiv\vartheta$ in $\degree$C) gives the following results (s. Table \ref{Hg Peak voltages and currents}):
\begin{table}[h!]
\centering
\begin{tabular}{c|c||c|c||c|c}
$U_\tecxt\text{1,max}^{175}\,/\,$V&$I_\tecxt\text{A,max}^{175}\,/\,$nA&$U_\tecxt\text{1,max}^{180}\,/\,$V&$I_\tecxt\text{A,max}^{180}\,/\,$nA&$U_\tecxt\text{1,max}^{185}\,/\,$V&$I_\tecxt\text{A,max}^{185}\,/\,$nA \\ \hline
12,233&	1,2567&	12,2&	1,07&	12,167&	0,9267 \\ \hline
16,9	&2,66&	16,8&	2,2433&	16,733&	1,9333 \\ \hline
21,467&	4,67&	21,367&	3,97&	21,367&	3,36 \\ \hline
26,2&	7,32&	26,167&	6,1833&	26,167&	5,2233 \\ \hline
31&	10,55&	30,933&	8,9233&	30,933&	7,54\\ \hline
35,8&	14,49&	35,7&	12,157&	35,633&	10,283\\ \hline
40,7&	18,95&	40,533&	15,953&	40,467&	13,483\\ \hline
45,6&	24,123&	45,4&	20,273&	45,2&	17,107\\ \hline
50,467&	29,947&	50,3&	25,17&	50,133&	21,237\\ \hline
55,4&	36,42&	55,167&	30,547&	54,867&	25,743
\end{tabular}
\caption{Mercury Peak voltages and currents for all three temperatures}
\label{Hg Peak voltages and currents}
\end{table}\\
The systematic uncertainty of $U$ is given by $\Delta U^\tecxt\text{sys}=\frac{1}{\sqrt{3}}\cdot 0,01\cdot U\,$ according to the instructions. The differences of the acceleration voltages are calculated as $U_\tecxt\text{diff}^\tecxt\text{max}=U_{i+1}^\tecxt\text{max}-U_i^\tecxt\text{max}$. For this, the results for the three different temperatures were combined. The statistic uncertainties of the individual peaks were caluclated via the usual formula. The statistic and systematic uncertainties of the difference voltages were calculated using the Gauss' law under the assumption of full correlation ($\Delta U_\tecxt\text{diff}=|\Delta U_{i+1}^\tecxt\text{max}-U_i^\tecxt\text{max}|$), s. Table \ref{Difference voltages and uncertainties for mercury}:
\begin{table}[h!]
\centering
\begin{tabular}{c|c|c}
$U_\tecxt\text{diff}\,/\,$V&$\Delta U_\tecxt\text{diff}^\tecxt\text{sys}\,/\,$V&$\Delta U_\tecxt\text{diff}^\tecxt\text{stat}\,/\,$V \\ \hline
4,6111&	0,0164&	0,0266 \\ \hline
4,5889&	0,0164&	0,0265\\ \hline
4,7778&	0,0307&	0,0276\\ \hline
4,7778&	0,0099&	0,0276\\ \hline
4,7556&	0,0294	&0,0275\\ \hline
4,8556&	0,0388	&0,0280\\ \hline
4,8333&	0,0709&	0,0279\\ \hline
4,9000&	0,0268&	0,0283\\ \hline
4,8444&	0,0525&	0,0280
\end{tabular}
\caption{Difference voltages and uncertainties for mercury}
\label{Difference voltages and uncertainties for mercury}
\end{table}\\
From this one can calulate the average excitation energy, according to $E_a=e\cdot\bra\langle {\Delta U_1^{\tecxt\text{max}}}\rangle$:
\begin{align*}
E_a^\tecxt\text{Hg}&=\underline{\kla\left( {4,772\pm 0,028\pm 0,032} \right)\,\tecxt\text{eV}}
\end{align*}
With $E_a=hf_a=\frac{hc}{\lambda_a}$, this gives the frequency and wavelength of the excitation ($c = 299762458\,\frac{\tecxt\text{m}}{\tecxt\text{s}}$, $e = 1.602176634 \cdot 10^{-19}\,$As and $h = 6.62607015 \cdot 10^{-34}\,$Js):
\begin{align*}
f_a^\tecxt\text{Hg}&=\underline{\kla\left( {1154\pm 7\pm 8} \right)\,\tecxt\text{THz}} \\
\lambda_a^\tecxt\text{Hg}&=\underline{\kla\left( {259,8\pm 1,5\pm 1,8} \right)\,\tecxt\text{nm}}
\end{align*}
Fig. \ref{Hg pair voltages per index} shows a plot of the excitation energies as a function of the pair index. Since the behaviour seems to be identical for all three temperatures, only one plot is shown:
\begin{figure}[h!]
\centering
\includegraphics[scale=1]{Hg_pair_voltages.png}
\caption{Hg-tube, Pair voltages per index with statistical error}
\label{Hg pair voltages per index}
\end{figure}\\
There seems to be a tendency towards larger difference voltages for higher indices. Considering the small uncertainties obtained from the data evaluation, this effect can hardly be explained as statistical fluctuation. \\
Due to higher acceleration voltages, the electrons pick up more energy $\delta$ along their mean free path between two collisions. Hence, the more energy the electrons have, the more collisions occur and the more energy they can pick up inbetween those. Denote $n$ as the number of a maximum in the Franck-Hertz curve. Then the electron energy should be weakly dependant of $n$ by $E_n=n\cdot( E_a+\delta)$.\\
With this in mind, the excitation should originate from the $6{}^3P_1\rightarrow 6{}^1S_0$ transition, which has an energy equivalent of $4,66\,$eV according to the instructions.
\subsection{Neon measurement}
Parameters:\\
Ne-tube, heating voltage $U_H=7,5\,V$, retarding voltage $U_2=2,0\,$V, control voltage $U_3=(1,5,\,2,\,2,5)\,$V. Each measurement was repeated twice to get a total of 9 curves. See Fig. \ref{Ne plot 1.5 V} for one of these as a general example.\\
\begin{figure}[h!]
\centering
\includegraphics[scale=1]{Ne_U3_2_5_plot.png}
\caption{Ne-tube, $U_3=1,5\,$V}
\label{Ne plot 1.5 V}
\end{figure}
Taking the average for the individual control voltages (superscript $\equiv U_3$ in V) gives the following results (s. Table \ref{Ne Peak voltages and currents}):
\begin{table}[h!]
\centering
\begin{tabular}{c|c||c|c||c|c}
$U_\tecxt\text{1,max}^{1,5}\,/\,$V&$I_\tecxt\text{A,max}^{1,5}\,/\,$nA&$U_\tecxt\text{1,max}^{2,0}\,/\,$V&$I_\tecxt\text{A,max}^{2,0}\,/\,$nA&$U_\tecxt\text{1,max}^{2,5}\,/\,$V&$I_\tecxt\text{A,max}^{2,5}\,/\,$nA \\ \hline
42,667&	1,180&		42,633&	1,460&		42,900&	1,830 \\ \hline
60,667&	1,627&		60,067&	2,250&		59,500&	3,020 \\ \hline
78,267&	2,253&		78,367&	3,333&		77,933&	4,730 \\ \hline
97,400&	3,027&		96,767&	4,740&		96,500&	6,897
\end{tabular}
\caption{Neon Peak voltages and currents for all three control voltages}
\label{Ne Peak voltages and currents}
\end{table}\\
The same procedure as for the Hg-measurement leads to the following difference voltages and uncertainties, s. Table \ref{Difference voltages and uncertainties for neon}:
\begin{table}[h!]
\centering
\begin{tabular}{c|c|c}
$U_\tecxt\text{diff}\,/\,$V&$\Delta U_\tecxt\text{diff}^\tecxt\text{sys}\,/\,$V&$\Delta U_\tecxt\text{diff}^\tecxt\text{stat}\,/\,$V \\ \hline
17,344&	0,247	&0,100 \\ \hline
18,111&	0,128&	0,105 \\ \hline
18,700&	0,019&	0,108
\end{tabular}
\caption{Difference voltages and uncertainties for neon}
\label{Difference voltages and uncertainties for neon}
\end{table}\\
With regard to the explaination of the increase of the difference voltages, the resulting transition energy, frequency and wavelength is calculated only from the first index pair to minimize said effect:
\begin{align*}
E_a^{\tecxt\text{Ne}}&=\underline{\kla\left( {17,34\pm 0,25\pm 0,10} \right)\,\tecxt\text{eV}} \\
f_a^\tecxt\text{Ne}&=\underline{\kla\left( {4194\pm 24\pm 60} \right)\,\tecxt\text{THz}} \\
\lambda_a^\tecxt\text{Ne}&=\underline{\kla\left( {71,5\pm 0,4\pm 1,0} \right)\,\tecxt\text{nm}}
\end{align*}
Fig. \ref{Ne pair voltages per index} shows a plot of the excitation energies as a function of the pair index:
\begin{figure}[h!]
\centering
\includegraphics[scale=1]{Ne_Pair_voltages.png}
\caption{Ne-tube, Pair voltages per index with statistical error}
\label{Ne pair voltages per index}
\end{figure}\\
The optical spectrum was aquired for a control voltage of $U_3=2,0\,$V. It yielded the following spectrum (s. Fig. \ref{Optical spectrum for neon}) and maxima with respective counts. The table also shows the measured and theoretical values of the energies corresponding to the transitions from the instructions (s. Table \ref{Peak wavelengths of the optical spectrum for neon}). I am unsure about the correct nomenclature for differentiating between the many different energy levels between the 18,3\,eV and 18,9\,eV line in the instructions. I tried to use the index for these states in a way to make it clear what transition I refer to graphically (?i? $\equiv$ i-th state from $\mathrm{2p^53p{}^3P_0}$ downwards). The energy levels were extraploated from the illustration:
\begin{figure}[h!]
\centering
\includegraphics[scale=1]{Ne_optical_spectrum.png}
\caption{Optical spectrum for neon}
\label{Optical spectrum for neon}
\end{figure}\\
\begin{table}[h!]
\centering
\begin{tabular}{c|c|c|c|c}
$\lambda_\tecxt\text{max}\,/\,$nm &Counts&$E_a^{\tecxt\text{meas}}\,/\,$eV&Transition&$E_a^\tecxt\text{theo}\,/\,$eV\\ \hline
586,0&364&2,116&$\mathrm{2p^53p{}^3P_0\rightarrow 2p^53s{}^1P_1}$&2,11  \\ \hline
614,7&165&2,017&$\mathrm{2p^53p{}^3P_{?5?}\rightarrow 2p^53s{}^3P_2}$&2,00  \\ \hline
640,5&172&1,936&$\mathrm{2p^53p{}^3P_{?6?}\rightarrow 2p^53s{}^3P_2}$&1,93  \\ \hline
650,8& 69&1,905&$\mathrm{2p^53p{}^3P_{?6?}\rightarrow 2p^53s{}^3P_1}$&1,89  \\ \hline
669,8& 54&1,851&$\mathrm{2p^53p{}^3P_{?4?}\rightarrow 2p^53s{}^1P_1}$&1,84  \\ \hline
704,0& 64&1,761&$\mathrm{2p^53p{}^3S_1\rightarrow 2p^53s{}^3P_2}$&1,73
\end{tabular}
\caption{Peak wavelengths and corresponding energies of the optical spectrum for neon}
\label{Peak wavelengths of the optical spectrum for neon}
\end{table}\\
The calculated transition energy is not in good agreement with the only allowed transition of atleast similar energy, which is the ${}^1P_1\rightarrow {}^1S_0$ transition with an energy equivalent of $16,79\,$eV according to the instructions.  
\subsection{Child's law for mercury}
Since $I_0\approx 0\,$nA, Child's law is $I_A\propto (U_1-U_K)^2$. Therefor, plotting $I_A$ versus $U_1^2$ and calculating the linear regression gives the result seen in Fig. \ref{Proving Childs law for Hg}. The dashed lines correspond to the systematical error of $U_1^2$, calculated from $\Delta U_1^2=2U_1\Delta U_1=\frac{1}{\sqrt{3}}\cdot 0,01\cdot 2U_1^2$. The statistical error is shown as x-errorbars:
\begin{figure}[h!]
\centering
\includegraphics[scale=1]{Childs_law_Hg.png}
\caption{Child's law for Hg with boundary curves}
\label{Proving Childs law for Hg}
\end{figure} 
Let $a$ denote the slope of the linear regression and $b$ its offset. The contact voltage $U_K$ can then be calulated from the zero-point of the linear regression, i.e. $U_K=\sqrt{\frac{-b}{a}}$. For the lack of a better idea, the statistical uncertainty was extrapolated linearly from the uncertainty of the smallest peak voltage, i.e. $\Delta U_M\approx \Delta U_1\cdot\frac{U_K}{U_{1,\tecxt\text{max}}^{185}}$. The systematic uncertainty of the contact voltage can be evaluated as follows:
\begin{align*}
\pm \Delta U_K^\tecxt\text{sys}&=U_K\frac{\pm\Delta a}{2a}=U_K\frac{a^\pm -a}{2a}
\end{align*}
Here $a^\pm$ refers to the slope of the boundary curves, ($a^+ =8,8666\cdot 10^{-3}\,\frac{\text{nA}}{\text{V}^2}$ and $a^-=8,5819\,\frac{\tecxt\text{nA}}{\tecxt\text{V}^2}$). The final result is:
\begin{align*}
U_K&=\underline{\kla\left( {8,40_{-0,05}^{+0,05}\pm 0,06} \right)\,\tecxt\text{V}}
\end{align*}
\section{Discussion}
With regard to the increased spacing between higher order voltage peaks, the final results for the Franck-Hertz curves are presented as a ratio of the first order maximum to the associated literature value. For mercury this gives:
\begin{align*}
\frac{E_a^\tecxt\text{Hg}}{E_a^\tecxt\text{Hg,lit.}}=\underline{\kla\left( {0,990\pm 0,006\pm 0,004} \right)}
\end{align*}
Despite the good agreement, the correct value of 1 is just barely within the uncertainty range, due to the small errors. Similarly for the neon result:
\begin{align*}
\frac{E_a^\tecxt\text{Ne}}{E_a^\tecxt\text{Ne,lit.}}=\underline{\kla\left( {1,033\pm 0,006\pm 0,015} \right)}
\end{align*}
Due to larger transition energies, the difference voltages between the neon peaks are higher. This also leads to broader peaks; they don't seem to be qualitatively different from mercury, just stretched out by a factor of about ${U_\tecxt\text{diff}^\tecxt\text{Ne}}/{U_\tecxt\text{diff}^\tecxt\text{Hg}}$.\\
Six peaks of the optical spectrum of neon could be identified energetically to atleast 1\% precision (except for the 704\,nm line with roughly 2\%). However as mentioned above, I am unsure about some of the corresponding transitions.\\ Since mercury does have very well-defined lines in the optical spectrum, I can not assign a good reason as to why it should not be possible to aquire said spectrum with the Hg-tube (according to the persistent lines listed in [\footnote{https://physics.nist.gov/PhysRefData/Handbook/Tables/mercurytable3.htm}]). \\
Child's law could be proven quite succesfully, as there was a very good agreement of the data to a $I_A\propto (U_1-U_K)^2$ dependance. The contact voltage was calulated to:
\begin{align*}
U_K&=\underline{\kla\left( {8,40_{-0,05}^{+0,05}\pm 0,06} \right)\,\tecxt\text{V}}
\end{align*}
\end{document}